\chapter{Gambling Addiction}

\section*{Definition and Overview}
Gambling addiction, also called gambling disorder or problem gambling, is a condition in which a person cannot control their gambling despite the harm it causes to themselves, their family, and their finances. It is recognised as a behavioural addiction, with features similar to substance addiction: craving, loss of control, tolerance (needing to bet more for the same thrill), withdrawal-like distress, and continued gambling despite consequences. Problem gambling affects a significant minority of people who gamble, and it is associated with financial ruin, relationship breakdown, mental illness, and suicide. It is treatable, and most people who seek help recover. This chapter explains the signs, causes, and treatments for gambling addiction.

\section*{Epidemiology}
Gambling is widespread in Australia, which has some of the highest rates of gambling in the world, including pokies (electronic gaming machines), sports betting, and lotteries. Around half of Australian adults gamble, and around 1% have a severe gambling problem, with a further 2--3% at moderate risk. Problem gambling affects more than 100,000 Australians and harms many more through the effects on families. It is more common in men, in people with low income, and in those with other addictions or mental illness. Online gambling has increased the accessibility and risk of harm.

\section*{Aetiology and Pathophysiology}
\begin{itemize}
    \item \textbf{The reward system:} gambling activates the brain's dopamine reward pathway, producing excitement that is reinforced by near-misses and intermittent wins
    \item \textbf{Variable reinforcement:} the unpredictable pattern of wins and losses is the most powerful driver of compulsive behaviour
    \item \textbf{The "chase":} after losses, the person bets more to recover them, deepening the losses
    \item \textbf{Cognitive distortions:} beliefs such as "I'm due for a win", luck rituals, and overestimating skill sustain gambling
    \item \textbf{Genetic and personality factors:} impulsivity and a family history of addiction increase risk
    \item \textbf{Environmental factors:} easy access, advertising, and gambling venues designed for continuous play
    \item \textbf{Co-occurring conditions:} depression, anxiety, alcohol use, and other addictions commonly coexist
    \item \textbf{The cycle:} gambling causes financial and relationship stress, which worsens mood, which drives further gambling
\end{itemize}

\section*{Clinical Features}
\begin{itemize}
    \item Spending more time and money gambling than intended, and difficulty stopping or cutting back
    \item Restlessness or irritability when trying to reduce gambling
    \item "Chasing" losses by gambling again soon after losing
    \item Lying to family and friends about gambling
    \item Borrowing money, selling possessions, or using savings and credit to gamble
    \item Jeopardising work, study, and relationships because of gambling
    \item Relying on others to relieve gambling-related financial problems
    \item Neglecting responsibilities, and gambling to escape stress or low mood
    \item Financial distress, debt, and in severe cases thoughts of suicide
\end{itemize}

\section*{Differential Diagnosis}
Gambling behaviour ranges from social to disordered, and other issues may coexist.
\begin{itemize}
    \item Social gambling, which is controlled, affordable, and not harmful
    \item Harmful gambling that has not yet reached the level of a disorder
    \item Bipolar disorder, in which excessive gambling can occur during manic episodes
    \item Other addictions, including alcohol and drugs, which often co-occur
    \item Depression and anxiety, which can both cause and result from gambling problems
    \item Attention deficit hyperactivity disorder, with impulsivity
    \item The severity and impact, not the amount gambled, define the disorder
\end{itemize}

\section*{When to Seek Medical Care}
Anyone who feels their gambling is out of control, or whose family is concerned, should seek help early. Doctors can assess, refer to specialist services, and treat co-occurring mental illness. Family members should also seek support.

\textbf{Call 000 (or local emergency services) for:}
\begin{itemize}
    \item Suicidal thoughts or behaviour, which are more common in gambling addiction
    \item Severe distress, panic, or collapse
    \item Any crisis involving family violence or harm
\end{itemize}

\textbf{Support lines:} Gambling Help Online (1800 858 858) and Lifeline (13 11 14) are available 24 hours a day.

\section*{Diagnosis and Investigations}
\begin{itemize}
    \item \textbf{Clinical interview:} the doctor explores gambling behaviour, its impact, and its effects on mood and relationships
    \item \textbf{Screening tools:} brief questionnaires such as the Problem Gambling Severity Index assess the level of harm
    \item \textbf{Mental health assessment:} screening for depression, anxiety, and other addictions
    \item \textbf{Financial and social assessment:} the extent of debt and family impact informs the treatment plan
    \item \textbf{Family involvement:} with consent, family perspectives are valuable
\end{itemize}

\section*{Management}
\begin{itemize}
    \item \textbf{Cognitive behavioural therapy:} the most effective treatment, addressing the thoughts and behaviours that drive gambling
    \item \textbf{Motivational interviewing:} builds the motivation to change
    \item \textbf{Counselling services:} free, confidential gambling counselling is available through services such as Gambler's Help and Relationships Australia
    \item \textbf{Financial counselling:} help with debt management and self-exclusion arrangements
    \item \textbf{Self-exclusion:} banning oneself from venues and online gambling sites
    \item \textbf{Medicines:} there is no medicine specifically approved for gambling disorder, but treating co-occurring depression, anxiety, and ADHD is important
    \item \textbf{Peer support:} groups such as Gamblers Anonymous provide community and accountability
    \item \textbf{Practical controls:} limiting cash and card access, closing betting accounts, and having a trusted person manage finances
    \item \textbf{Family support:} counselling for partners and children affected by the gambling
\end{itemize}

\section*{Complications}
\begin{itemize}
    \item Severe debt, bankruptcy, and loss of housing
    \item Relationship breakdown and family conflict
    \item Job loss and reduced work performance
    \item Depression, anxiety, and substance use
    \item Suicide, which is significantly more common in gambling disorder
    \item The intergenerational impact on children
\end{itemize}

\section*{Prognosis and Outcomes}
Gambling addiction is treatable, and most people who seek help make significant and lasting progress. Cognitive behavioural therapy and counselling produce good outcomes, and combining treatment with practical measures such as self-exclusion and financial management greatly improves the chance of recovery. Relapse can occur, particularly in times of stress, but people who maintain support and controls recover over time. Families also recover with support. The most important step is asking for help, and free services make treatment accessible to everyone.

\section*{Prevention and Self-Care}
\begin{itemize}
    \item Set a limit on money and time before gambling, and never gamble to recover losses
    \item Avoid gambling when stressed, low, or affected by alcohol
    \item Use venue and online self-exclusion if gambling is hard to control
    \item Keep gambling separate from credit, savings, and essential money
    \item Talk to a doctor or call the gambling helpline early, before the harm grows
    \item Seek treatment for co-occurring depression, anxiety, or alcohol problems
    \item Support a partner or family member by encouraging help and attending counselling
    \item Never gamble to escape emotional distress
\end{itemize}

\section*{Sources and Further Reading}
The following reputable sources provide further information for healthcare students and patients seeking reliable, current advice about gambling addiction.
\begin{itemize}
    \item Gambling Help Online. \textit{Support and counselling.} https://www.gamblinghelponline.org.au/
    \item Gambler's Help. \textit{Free counselling services.} https://www.gamblershelp.com.au/
    \item Healthdirect Australia. \textit{Problem gambling.} https://www.healthdirect.gov.au/problem-gambling
    \item Relationships Australia. \textit{Support for gambling and families.} https://www.relationships.org.au/
    \item NHS. \textit{Problem gambling.} https://www.nhs.uk/conditions/problem-gambling/
    \item World Health Organization. \textit{Gambling disorder.} https://www.who.int/
\end{itemize}
